\documentclass[conference]{IEEEtran}
\usepackage{amsmath,amssymb}
\begin{document}
\subsection{Geospatial Routing and Hydrological Catchment Gating}
To prevent cross-watershed contamination, incoming sensor feeds undergo topological spatial routing. Each observation is annotated with coordinates and assigned to a hydrological sub-basin via USGS HUC-10 boundaries. Let (x,\ell) \equiv \mathbf{1}[\mathrm{HUC}_{10}(x) = \mathrm{HUC}_{10}(\ell)]$ denote sub-basin concordance. The routing predicate is:
\begin{equation}
\label{eq:geo_gate}
g(x, \ell) =
\begin{cases}
\mathcal{Z}_{\mathrm{dir}}, & d(x,\ell) \le 5\,\text{km} \land m(x,\ell),\\[2pt]
\mathcal{Z}_{\mathrm{catch}}, & d(x,\ell) \le 25\,\text{km} \land m(x,\ell),\\[2pt]
\mathcal{Z}_{\mathrm{reg}}, & 25 < d(x,\ell) \le 60\,\text{km} \land x \in \mathcal{D}_{\mathrm{meteo}},\\[2pt]
\mathcal{Z}_{\mathrm{rej}}, & \text{otherwise,}
\end{cases}
\end{equation}
where (\cdot)$ is the Haversine great-circle distance, $\mathcal{Z}_{\mathrm{dir}}$ is direct local, $\mathcal{Z}_{\mathrm{catch}}$ is catchment corridor, $\mathcal{Z}_{\mathrm{reg}}$ is regional meteorology, and $\mathcal{Z}_{\mathrm{rej}}$ is rejected out-of-bounds.
\end{document}
